\section{Discussion}\label{sec:discussion}

\subsection{Answer to the question}
\begin{itemize}[itemsep=1pt]
\item \textbf{Mechanism.} Fiber orientation controls in-plane elastic waves by
  \emph{channeling} their energy along the fibers. The control is the redirection of a
  plate's vibration energy toward a target, by a few times ($\WAbsBand$--$\WRedirBandMean\times$
  band-averaged for the lens, depending on the absorbing layer), over a broad band.
\item \textbf{What bounds it.} The lamina anisotropy (none for isotropic material, growing
  with the channeling share) and the tow curvature (no loss while the minimum tow radius is below about $0.7$
  along-fiber wavelengths, then a gradual decline of the best design found).
\item \textbf{Cancellation.} Canceling the field of a finite scatterer (a full cloak) has to
  rely on interference, because exact cloaking needs anisotropy that no lamina has, and
  strong cancellation holds only near the design frequency (the curvature-limited design
  keeps a moderate $\WCloakCurvBandMin$--$\WCloakCurv\times$ across $\omega\in[40,60]$); the inextensible-fiber limit \cite{olsson2011} and
  ground-plane cloaks \cite{chen2011,popa2011} escape this bound.
\end{itemize}

\begin{table}[h]
\centering\footnotesize
\caption{Status of the predictions.}\label{tab:verdictW}
\begin{tabular}{@{}p{0.30\linewidth}p{0.36\linewidth}p{0.28\linewidth}@{}}
\toprule
prediction & evidence & verdict\\ \midrule
W1 isotropic plate: no control (identity; consistency check) & Remark~\ref{rem:iso}; Table~\ref{tab:dose} ($s=0$) & holds (exact)\\
W2 redirection is broadband, extra gain is resonant & Table~\ref{tab:lensband}, Fig.~\ref{fig:lensband}; damping test (App.~\ref{app:verify}) & holds for $\omega\ge37.5$ ($G_\varphi=2.3$ at $35$); fails at long wavelength ($G_\varphi<1$ at $\omega=20$ and $25$)\\
W3 energy flux follows the fibers & lens: $\WFluxDevBandMin$--$\WFluxDevBandMax^\circ$ across the band (Fig.~\ref{fig:flux}); unoptimized controls: $\WCtrlDevMin$--$\WCtrlDevMax^\circ$ & holds\\
W4 redirection grows with anisotropy & Table~\ref{tab:dose}, Fig.~\ref{fig:limitsW}a & holds coarsely (1 at isotropy, several times once $E_{f1}/E_{f2}\gtrsim3$); finer steps are within optimizer variability\\
W5 redirection is unaffected while $R\lesssim0.7\lambda$, then declines; strong cloaking is narrowband & Tables~\ref{tab:front},~\ref{tab:cloakband}; multi-start front & partly holds: redirection is bounded by the tow radius, but declines gradually rather than failing at $R\sim\lambda$ (front values are lower bounds)\\
\bottomrule
\end{tabular}
\end{table}

\subsection{Consequences for how wave-control designs are reported}
A finite, lightly damped plate driven at one frequency is a resonator. An optimizer
maximizing a field energy at that frequency exploits the resonance whenever it can. The
lens shows the size of the effect: $\WGainForty\times$ at the design frequency, of which
$\WGrowthForty\times$ is field growth, and $\WGainMinusFive\times$ only $6\%$ away. We
therefore reaffirm three practices, established for wave topology optimization
\cite{jensen2005,andkjaer2013}, applied in this paper, and needed before single-frequency
gains are compared across studies:
\begin{enumerate}[itemsep=1pt]
\item report redirection (a share of energy) alongside, or instead of, raw gains;
\item evaluate every design across a band;
\item optimize band-averaged objectives when the application has a band
  (Sec.~\ref{sec:panel}).
\end{enumerate}

\subsection{Limitations and scope}
\begin{itemize}[itemsep=2pt]
\item \textbf{Model.} The model is a single plane-stress lamina, with no through-thickness
  or flexural waves, no tow gaps or overlaps, and constant thickness.
\item \textbf{Absorbing layer.} The sponge reflects $|R|\approx\WReflMin$--$\WReflMax$ of the
  along-fiber qP wave, so finite-plate standing waves remain. With a sponge twice as wide
  ($|R|\le\WReflThickMax$) the resonant growth of the lens persists ($\WThickGrForty\times$ at
  $\omega=40$), but the redirection is lower: $\WThickPForty\times$ at $\omega=40$ and
  $\WThickBand\times$ band-averaged, against $\WRedirForty\times$ and $\WThinBand\times$
  (Appendix~\ref{app:verify}). This is a systematic bias, not scatter: reflections from the layer raise the
  target share. Re-optimizing the lens with the better absorber does not recover the
  difference (band-averaged $\WAbsBand\times$, against $\WAbsBandOld\times$ for the
  original design with that absorber, and $\WAbsG\times=\WAbsP\times$ redirection
  $\times\,\WAbsGr\times$ growth at $\omega=40$). Absolute redirection values are therefore
  about a quarter lower with a near-non-reflecting boundary; the conclusions, which rest
  on the split between broadband redirection and resonant growth, do not change.
\item \textbf{Discretization.} Lens metrics change by at most about $7\%$ under $2\times$ mesh
  refinement. The panel's sensor and payload ratios change by at most $\WPanMeshSenDev\%$ on a
  finer mesh, its rim ratios by up to $\WPanMeshRimDevAbs$ in absolute terms. The optimized cloak is mesh-tuned ($36\times\to\WCloakFineEval\times$ on a
  finer mesh), but the capability persists ($\WCloakFineReopt\times$ re-optimized there);
  we checked two meshes, not a converged sequence.
\item \textbf{Optimality.} The optimizations are local (MMA, from straight fibers unless
  stated; the front also uses warm and random starts). The dose and front studies
  re-optimize each point independently, so their trends carry optimizer noise, and the
  front values are lower bounds on what each limit allows.
\item \textbf{Ray picture.} Channeling is argued from plane-wave theory and confirmed by the
  flux field. A full ray-tracing prediction of the redirection is left for future work.
\end{itemize}

\subsection{Outlook}
The channeling picture suggests where orientation design should go next:
\begin{itemize}[itemsep=1pt]
\item routing problems around prescribed holes and joints, where tows must wrap
  obstacles within the curvature limit;
\item multi-objective layouts that trade wave routing against stiffness;
\item combined topology and toolpath design, in which voids are cut rather than softened
  so that energy cannot pool in compliant regions;
\item a ray-theoretic predictor of redirection from the fiber field alone.
\end{itemize}

\section{Conclusions}\label{sec:conclusion}
\begin{enumerate}[itemsep=2pt]
\item In a CFRP lamina, wave energy is skewed toward the fibers by up to $65^\circ$, and
  $88\%$ of quasi-P wave normals carry their energy within $20^\circ$ of the fiber
  direction (Sec.~\ref{sec:theory}).
\item An orientation-only lens redirects energy into its target by $\WRedirForty\times$ at
  its design frequency and by at least $\WRedirBandMin\times$ over $\omega\in[37.5,60]$;
  with a better absorbing layer, re-optimized, by $\WAbsP\times$ and at least
  $\WAbsMinFortySixty\times$ over $[40,60]$.
  Its $\WGainForty\times$ single-frequency gain is mostly a resonant growth of the whole
  field.
\item In the optimized lens, the energy flux follows the local fibers: mean deviation
  $\WFluxDev^\circ$ at the design frequency and $\WFluxDevBandMin$--$\WFluxDevBandMax^\circ$
  across the band, and the same holds in curved fields that were never optimized.
\item Redirection vanishes exactly for isotropic material and grows with anisotropy. It
  is unaffected while the minimum tow radius is at most $0.7\lambda$; beyond it the best
  design found falls to about half by $R\approx1.2\lambda$ and to $\WFrontTight\times$ by
  $R\approx2.3\lambda$ (best of $\WFrontNStart$ starts per limit, a lower bound).
\item A bounded-anisotropy cloak reduces the far-field mismatch $\WCloakPeak\times$ at its
  design frequency, but only $\WCloakMinus$--$\WCloakPlus\times$ at $\pm5\%$; under the
  curvature limit it holds a moderate $\WCloakCurvBandMin$--$\WCloakCurv\times$ over
  $\omega\in[40,60]$. Transformation directions alone give
  $\WCloakAnalytic\times$.
\item Band-averaged design delivers multifunctional fiber layouts whose performance holds
  across the operating band.
\end{enumerate}

\paragraph{Data and code availability.} The solver, result files and the scripts that generate
the reported results, tables and figures (Appendix~\ref{app:verify} and the parameter tables) are in the
repository \url{https://github.com/jtempelman95/wavetopo} (branch \texttt{paper-revision}; see \texttt{docs/paper/README\_REPRODUCE.md}).

\paragraph{Competing interests.} The author declares no competing interests.
